% MT022 / explanation — canonical editable content.
任务有50,000名客户、25个特征，需要预测购买类别，同时分析哪些特征有帮助。按配套答案，可选择随机森林（random forest，RF）。


第一个理由是\textbf{能处理非线性和特征交互}。例如“收入高”未必单独意味着购买，可能还取决于年龄与使用频率。树在不同分支用不同条件判断，多棵树投票可以表示这类关系，不必预先手写所有交互项。


第二个理由是\textbf{能辅助分析特征的预测作用}。例如打乱某个特征后，验证表现明显下降，说明模型依赖它提供的信息，这叫置换重要性（permutation importance）。它能帮助找业务线索，但说明的是预测关联；不能由此断言改变收入就一定改变购买决定。


另外两种模型各写一个与本题有关的不足即可：



\begin{itemize}

\item \textbf{普通逻辑回归：} 在给定25个特征上使用线性分数，可能遗漏复杂交互。它其实很便于解释，也可以扩展特征，所以不能写成“LR没有解释性”。

\item \textbf{常规核SVM：} 需要处理样本之间的核值。完整核矩阵有\(50{,}000^2=2.5\times10^9\)个元素；若每个用8字节，单是该矩阵就约20 GB（十进制）。这是说明规模压力的计算，不是断言所有实现都必须保存整个矩阵。

\end{itemize}


实际谁更准仍要验证；原题允许有依据的模型选择。


\textbf{出处：} 2021A期中 Q9，原卷第3页（10分）


\textbf{答案依据：} 2021A期中 Q9 随卷／配套参考答案；条件与解释经复核。
